\documentclass[pad]{kiritex}

\title{高等数学习题集}
\author{风待桐叶}
\date{2026 年 7 月}

\begin{document}

\maketitle

\tableofcontents
\newpage

\section{极限与连续}

\subsection{极限的计算}

\begin{problem}[基础]
  求极限 $\displaystyle\lim_{x\to 0}\frac{\sin x}{x}$。
  \answer{极限值为 $1$。}
\end{problem}

\begin{problem}[基础]
  求极限 $\displaystyle\lim_{x\to\infty}\left(1+\frac{1}{x}\right)^x$。
  \answer{极限值为 $\mathrm{e}$。}
\end{problem}

\begin{problem}[综合]
  计算下列极限：
  \begin{problemsubparts}
    \item $\displaystyle\lim_{x\to 0}\frac{\tan x-\sin x}{x^3}$
    \item $\displaystyle\lim_{x\to 0}\frac{1-\cos x}{x^2}$
  \end{problemsubparts}
  \answer{分别为 $\frac12$ 与 $\frac12$。}
\end{problem}

\subsection{重要极限}

\begin{problem}[填空]
  重要极限 $\displaystyle\lim_{x\to 0}\frac{\sin x}{x}$ 的值是 \problemblock*
  \answer{$1$。}
\end{problem}

\begin{problem}[提高]
  求极限 $\displaystyle\lim_{x\to 0}(1+2x)^{\frac{1}{x}}$。
  \answer{极限值为 $\mathrm{e}^{2}$。}
\end{problem}

\subsection{连续性}

\begin{problemtf}[判断]
  若函数在一点可导，则它在该点连续。
  \answer{正确。}
\end{problemtf}

\begin{problem}[讨论]
  设 $f(x)=\begin{cases}
    x^2\sin\frac{1}{x}, & x\neq 0\\
    0, & x=0
  \end{cases}$，讨论 $f(x)$ 在 $x=0$ 处的连续性与可导性。
  \answer{$f(x)$ 在 $0$ 处连续且可导，且 $f'(0)=0$。}
\end{problem}

\begin{problemchoice}[选择]
  极限 $\displaystyle\lim_{x\to 0}\frac{\sin x}{x}$ 的值是
  \choice{0}
  \choice{1}
  \choice{$\infty$}
  \choice{不存在}
  \answer{B。}
\end{problemchoice}

\section{导数与微分}

\subsection{基本求导法则}

\begin{problem}[基础]
  求 $y=x^3-3x^2+2x-1$ 的导数。
  \answer{$y'=3x^2-6x+2$。}
\end{problem}

\begin{problem}[基础]
  求 $y=\ln(\sin x)$ 的导数。
  \answer{$y'=\cot x$。}
\end{problem}

\subsection{高阶导数}

\begin{problem}[综合]
  设 $y=\mathrm{e}^{x}\cos x$，求 $y''$。
  \answer{$y''=-2\mathrm{e}^{x}\sin x$。}
\end{problem}

\begin{problem}[提高]
  设 $y=\dfrac{1}{1-x}$，求 $y^{(n)}$。
  \answer{$y^{(n)}=\dfrac{n!}{(1-x)^{n+1}}$。}
\end{problem}

\subsection{导数的应用}

\begin{problem}[综合]
  求曲线 $y=x^3-3x$ 的拐点与凹凸区间。
  \answer{拐点为 $(0,0)$；$(-\infty,0)$ 上凹，$(0,\infty)$ 上凸。}
\end{problem}

\begin{problem}[综合]
  求函数 $f(x)=x^3-3x^2-9x+5$ 在 $[-2,4]$ 上的最大值与最小值。
  \answer{比较端点与驻点 $x=-1,3$ 的函数值即可。}
\end{problem}

\subsection{微分中值定理}

\begin{problem}[证明]
  设 $f(x)$ 在 $[0,1]$ 上连续，$(0,1)$ 内可导，且 $f(0)=0$，$f(1)=1$。
  证明：存在 $\xi\in(0,1)$ 使得 $f'(\xi)=2\xi$。
  \answer{对 $F(x)=f(x)-x^2$ 使用 Rolle 定理。}
\end{problem}

\begin{problem}[证明]
  验证函数 $f(x)=\ln x$ 在 $[1,\mathrm{e}]$ 上满足 Lagrange 中值定理的条件，并求 $\xi$。
  \answer{$1/\xi=(\ln\mathrm{e}-\ln1)/(\mathrm{e}-1)$，故 $\xi=\mathrm{e}-1$。}
\end{problem}

\section{不定积分}

\subsection{分部积分法}

\begin{problem}[基础]
  计算 $\displaystyle\int x\,\mathrm{e}^{x}\,\mathrm{d}x$。
  \answer{$(x-1)\mathrm{e}^{x}+C$。}
\end{problem}

\begin{problem}[基础]
  计算 $\displaystyle\int \ln x\,\mathrm{d}x$。
  \answer{$x\ln x-x+C$。}
\end{problem}

\subsection{有理函数积分}

\begin{problem}[综合]
  计算 $\displaystyle\int \frac{1}{x^2-1}\,\mathrm{d}x$。
  \answer{$\frac12\ln\left|\frac{x-1}{x+1}\right|+C$。}
\end{problem}

\begin{problem}[综合]
  计算 $\displaystyle\int \frac{x+1}{x^2+4x+3}\,\mathrm{d}x$。
  \answer{先分解 $x^2+4x+3=(x+1)(x+3)$。}
\end{problem}

\subsection{三角函数的积分}

\begin{problem}[基础]
  计算 $\displaystyle\int \sin^2 x\,\mathrm{d}x$。
  \answer{$\frac{x}{2}-\frac{\sin 2x}{4}+C$。}
\end{problem}

\begin{problem}[综合]
  计算 $\displaystyle\int \frac{1}{1+\cos x}\,\mathrm{d}x$。
  \answer{$\tan\frac{x}{2}+C$。}
\end{problem}

\subsection{换元法}

\begin{problem}[基础]
  计算 $\displaystyle\int \frac{\ln x}{x}\,\mathrm{d}x$。
  \answer{$\frac12(\ln x)^2+C$。}
\end{problem}

\begin{problem}[基础]
  计算 $\displaystyle\int \frac{1}{\sqrt{1-x^2}}\,\mathrm{d}x$。
  \answer{$\arcsin x+C$。}
\end{problem}

\end{document}
